\chapter{Introduction}

\section{Background and motivation}

Online bipartite matching is the canonical model of irrevocable
sequential allocation. Resources are known in advance; requests arrive
one at a time; each must be matched to an available compatible resource
(or dropped) immediately, before the rest of the input is seen. It is
the abstraction behind online advertising, ride-hailing dispatch, organ
exchange, and request routing in modern serving systems. Because
decisions cannot be undone, the difficulty lies entirely in committing
well under uncertainty about the future.

A recent and influential idea is to give such an algorithm a
\textbf{prediction} about the input (which resources will be contended,
or how many requests of each kind will arrive), and to ask for two
guarantees at once: \textbf{consistency}, meaning near-optimal
performance when the prediction is good, and \textbf{robustness},
meaning no worse than a prediction-free algorithm when the prediction is
wrong. (We call the object a \emph{prediction} throughout, and
\emph{advice} only where the question is whether to trust it; the two
denote the same thing.) For online matching specifically, two families
of such \emph{learning-augmented} algorithms have appeared:
MinPredictedDegree, which consumes a prediction of each resource's
degree (how many kinds of request will compete for it, i.e.~how
contended it will be), and a family of \emph{test-and-fallback} schemes,
which test a request-type prediction on a short prefix of arrivals
before deciding whether to trust it.

These algorithms come with attractive worst-case guarantees, yet, as
this thesis demonstrates experimentally, their average-case behavior is
strikingly muted: on realistic workloads the sophisticated predictions
rarely \emph{help}, while the algorithms' real value is that they do not
\emph{hurt}. This gap between the worst-case promise and the
average-case reality is the puzzle that motivates this thesis.

Moreover, the published work we examined studied these algorithms
separately: each on its own input model and its own notion of prediction
error, and largely through theory. There was therefore no common
empirical basis on which to compare them or quantify how much a
prediction buys. This thesis provides that basis and investigates the
gap.

\section{Research objectives}

Performance is measured throughout as the ratio of the matching an
algorithm produces to the offline optimum on the same instance (§3.1
makes this precise). The thesis is organized around three questions:

\begin{enumerate}
\def\labelenumi{\arabic{enumi}.}
\tightlist
\item
  \textbf{How do the learning-augmented online-matching algorithms
  actually compare}, head to head, under a single prediction-error model
  and on common inputs, and how much does a prediction help, at what
  cost, on realistic data?
\item
  \textbf{What are the failure modes of the adaptive test-and-fallback
  mechanism}: the cost of its test, the calibration of its accept/reject
  decision, and where it goes wrong?
\item
  \textbf{Why does the average-case experience differ so sharply from
  the worst-case promise}: is the observed wall an artifact of
  particular algorithms and generators, or is it necessary?
\end{enumerate}

The first two are answered experimentally. The third receives a bounded
theoretical treatment---one proved trade-off inequality and a
quantitative interpretation---but is not closed.

\section{Contributions}

\begin{itemize}
\item
  \textbf{A unified experimental benchmark} (Chapter 4; Q1): a
  head-to-head comparison of these families, on one experimental
  harness: a single code path in which every algorithm sees the same
  graphs, predictions and optimum, with a \emph{structured} error model
  (errors follow the instance's structure, not i.i.d. noise) and
  confidence intervals.
\item
  \textbf{A characterization of predictions as robustness insurance}
  (Chapters 4, 7; Q1): unguarded following crashes below the advice-free
  baseline, two mechanisms (structural and adaptive) restore safety with
  different trade-offs, and the consistency upside is small, so the
  value is downside protection. Validated on synthetic graphs, six
  real-world graphs and real traces, including with a cheap, non-ML
  historical predictor.
\item
  \textbf{An empirical engagement with the order-error theory} (Chapter
  5; Q1): MinPredictedDegree's loss is governed by a Kendall-\(\tau\)
  order error, the fraction of resource pairs ranked in the wrong
  relative order, onto which several error models collapse.
  Order-dependence itself is Aamand--Chen--Indyk's result; ours is the
  characterization of their bound's looseness and saturation, and of the
  right order measure.
\item
  \textbf{An empirical study of test-and-fallback} (Chapter 6; Q2): its
  testing cost, a threshold-calibration pathology in which a \emph{more
  accurate} test yields a \emph{worse} decision, the resolution limit
  that recalibration exposes, and an irrevocability penalty that
  explains why matching requires \emph{test-then-commit}.
\item
  \textbf{A quantified account of why the wall stands} (Chapter 6; Q3).
  By the \emph{wall} we mean that on average-case inputs neither a
  better algorithm nor a better prediction moves the ratio appreciably.
  We trace it to a \emph{resolution limit} (no practical threshold can
  capture an upside below its own estimator's noise) that persists
  across the whole difficulty range (Figure 6.4).
\end{itemize}

Alongside these, the thesis documents (Chapter 8) the exploratory
directions it pursued and the negative results they returned: an attempt
to learn the predictor for extra performance, and an attempt to find a
serving regime where predictions genuinely help. Both reinforce the
central finding, and the question they raise --- is the wall forced? ---
is taken up in the concluding outlook (§10.2). That outlook is a
discussion rather than a further contribution: it proves one
consistency/robustness trade-off inequality and then reads our own
experiments through it. No theorem beyond that inequality is claimed
anywhere in this thesis.

\section{Thesis organization}

The chapters group by the three questions above. Chapters 2 and 3 build
the common ground: the background (input models, the learning-augmented
paradigm, the specific algorithms, and the distribution-testing results
behind their tests) and then our own model, methodology and harness,
validated by reproducing a subset of the Borodin et al.~study. Chapters
4 to 7 answer Q1 and Q2: the unified benchmark and its four findings
(4), what governs the small loss (5), the test-and-fallback mechanism in
depth (6), and external validity on real predictors and real graphs (7).
Chapters 8 to 10 take up Q3 and the boundaries: the exploratory
directions and their negative results (8), the AI-inference serving case
study (9), and a conclusion that summarizes, gives a theoretical outlook
on why the wall stands (§10.2), evaluates the project critically against
these objectives (§10.3), and states limitations and future work (10).
One sentence recurs throughout, and §10.1 states it in full: \textbf{on
average-case online matching, predictions are robustness insurance
rather than a performance lever.}
